% Part: incompleteness
% Chapter: theories-computability
% Section: consis-with-q

\documentclass[../../../include/open-logic-section]{subfiles}

\begin{document}

\olfileid{inc}{tcp}{con}

\olsection{$\Th{Q}$ शी सुसंगत उपपत्ती अनिर्णेय असतात}

पुढील प्रमेय असे सांगते की केवळ $\Th{Q}$ अनिर्णेय आहे असे नाही,
तर $\Th{Q}$ शी विसंगती निर्माण न करणारी कोणतीही उपपत्ती अनिर्णेय असते.

\begin{thm}
$\Th{T}$ ही अंकगणिताच्या भाषेतील $\Th{Q}$ शी सुसंगत असलेली
कोणतीही उपपत्ती असू द्या (म्हणजे $\Th{T} \cup \Th{Q}$ सुसंगत आहे).
मग $\Th{T}$ अनिर्णेय असते.
\end{thm}

\begin{proof}
$\Th{Q}$ ची $!Q_1$, \dots, $!Q_8$ ही सांत स्वयंसिद्धके आहेत,
हे आठवा. त्यांच्या जागी $!E = !Q_1 \land
\dots \land !Q_8$ हे एकच स्वयंसिद्धकही घेता येते.

$\Th{T}$ ही $\Th{Q}$ शी सुसंगत असलेली निर्णेय उपपत्ती आहे
असे समजू. पुढील संच घेऊ:
\[
C = \Setabs{!A}{\Th{T} \Proves !E \lif !A}.
\]
निश्चित सूत्राशी अभिव्यंजन तयार करणे संगणनक्षम आहे;
गृहीत उपपत्ती निर्णेय असल्याने या संचातील सदस्यत्वही निर्णेय आहे.
$C$ हा $\Th{Q}$ आणि $\Th{\bar Q}$ यांना संगणनक्षमपणे
विलग करणारा संच ठरेल, असे दाखवू; त्यामुळे व्याघात मिळेल.
प्रथम, $!A$ हे $\Th{Q}$ मध्ये असेल, तर $\Th{Q}$ च्या
स्वयंसिद्धकांपासून $!A$ सिद्ध होते. निगमन प्रमेयानुसार
प्रथम-क्रम तर्कशास्त्रात $!E \lif !A$ ची !!a{derivation}
असते. तार्किक प्रमेय असल्याने ते गृहीत निर्णेय उपपत्तीतही सिद्ध होते.
म्हणून $!A$ हे $C$ मध्ये आहे.

दुसरीकडे, $!A$ हे $\Th{\bar Q}$ मध्ये असेल, तर प्रथम-क्रम
तर्कशास्त्रात $!E \lif \lnot !A$ ची सिद्धता असते.
$\Th{T}$ हे $!E \lif !A$ सुद्धा सिद्ध करत असेल, तर
$\Th{T}$ हे $\lnot !E$ सिद्ध करेल; अशा वेळी
$\Th{T} \cup \Th{Q}$ विसंगत होईल. पण
$\Th{T} \cup \Th{Q}$ सुसंगत आहे असे आपण गृहीत धरले आहे.
म्हणून $\Th{T}$ हे $!E \lif !A$ सिद्ध करत नाही आणि $!A$
हे $C$ मध्ये नाही.

म्हणजे $!A$ हे $\Th{Q}$ मध्ये असेल तर $C$ मध्ये असते,
आणि $!A$ हे $\Th{\bar Q}$ मध्ये असेल तर $\Complement{C}$
मध्ये असते, हे दाखवले. त्यामुळे $C$ हा संगणनक्षम विलगक
ठरतो; हाच अपेक्षित व्याघात आहे.
\end{proof}

हे प्रमेय अतिशय प्रभावी आहे. उदाहरणार्थ, त्यावरून पुढील निष्कर्ष मिळतो:

\begin{cor}
  अंकगणिताच्या भाषेसाठी प्रथम-क्रम तर्कशास्त्र (म्हणजे
  $\Setabs{!A}{\text{$!A$ हे प्रथम-क्रम तर्कशास्त्रात सिद्ध होते}}$
  हा संच) अनिर्णेय आहे.
\end{cor}

\begin{proof}
प्रथम-क्रम तर्कशास्त्र हा $\emptyset$ पासून निष्पन्न होणाऱ्या
वाक्यांचा संच आहे; आणि तो रिक्त स्वयंसिद्धकसंच $\Th{Q}$ शी सुसंगत आहे.
\end{proof}

\end{document}
